% SPDX-License-Identifier: CC-BY-SA-4.0
% Author: Matthieu Perrin
% Part: <Nom de la partie>
% Section: <Nom de la section>
% Sub-section: <Nom de la sous-section>  % (facultatif, laisser vide si non utilisé)
% Frame: <Titre de la slide>

\begingroup

\newcommand\sep{\,\mathord{\circ}\,}

\begin{frame}{Indécidabilité de PCP}

  \vspace{-2mm}
  \begin{block}{Théorème}
    Le problème de correspondance de Post (\textsc{pcp}) est indécidable. 
  \end{block}

  \vspace{1mm}
  \begin{block}{Démonstration}
    \begin{itemize}
    \item On propose la réduction $\varphi$ de $\textsc{pcpi}$ sur $\Sigma$ vers $\textsc{pcp}$ sur $\Sigma\cup \{\circ\}$ :
    \end{itemize}
    \vspace{2mm}

    $$
    \varphi : \left\{\begin{array}{r@{~~}c@{~~}l}
    \mathcal{I}(\textsc{pcpi})
    &\rightarrow&
    \mathcal{I}(\textsc{pcp})
    %
    \vspace{2mm}\\
    %
    \left\langle\, D, \VDomino[1.3]{\gamma_1\ldots \gamma_{|\gamma|}}{\delta_1\ldots \delta_{|\delta|}} \,\right\rangle
    &\mapsto&
    \left\{\,
    \VDomino[2.2]{\sep \alpha_1 \sep \ldots \sep \alpha_{|\alpha|} \phantom{\sep}}{\phantom{\sep} \beta_1 \sep \ldots \sep \beta_{|\beta|} \sep}
    ~\middle|~\VDomino[1.3]{\alpha_1 \ldots \alpha_{|\alpha|}}{\beta_1 \ldots \beta_{|\beta|}} \in D
    \,\right\}
    \vspace{1mm}\\
    &&\cup
    \left\{\,
    \VDomino[2.2]{\sep \gamma_1 \sep \dots \sep \gamma_{|\gamma|} \phantom{\sep}}{\sep \delta_1 \sep \ldots \sep \delta_{|\delta|} \sep},
    \VDomino{\sep\sep}{\phantom{\sep}\sep}
    \,\right\}
    \end{array}\right.
    $$

    \vspace{3mm}
    \begin{itemize}
    \item $\varphi$ est bien calculable
    \item $\forall x \in \mathcal{I}(\textsc{pcpi}), \quad \varphi(x) \in \textsc{pos}(\textsc{pcp}) \quad \Leftrightarrow \quad x \in \textsc{pos}(\textsc{pcpi})$
    \item Or \textsc{pcpi} est indécidable, donc \textsc{pcp} est indécidable
    \end{itemize}

  \end{block}

\end{frame}

\endgroup
